% ===========================================================================
%  DSCSP: Diffusion-Stego with Conditional Sampling Paths
%  for Robust Image Steganography — IEEE Conference Format
% ===========================================================================
\documentclass[conference]{IEEEtran}

% ─── Packages ──────────────────────────────────────────────────────────────
\usepackage{amsmath, amssymb, amsfonts}
\usepackage{graphicx}
\usepackage{booktabs}
\usepackage{hyperref}
\usepackage{algorithm}
\usepackage{algorithmic}
\usepackage{multirow}
\usepackage{enumitem}
\usepackage{caption}
\usepackage{subcaption}
\usepackage{xcolor}
\usepackage{bm}
\usepackage{microtype}
\usepackage{cite}

\hypersetup{
    colorlinks=true,
    linkcolor=blue!70!black,
    citecolor=green!50!black,
    urlcolor=blue!60!black,
}

% ─── Math shortcuts ────────────────────────────────────────────────────────
\newcommand{\R}{\mathbb{R}}
\newcommand{\E}{\mathbb{E}}
\newcommand{\Lcal}{\mathcal{L}}
\newcommand{\Acal}{\mathcal{A}}
\newcommand{\Ncal}{\mathcal{N}}
\newcommand{\Ucal}{\mathcal{U}}
\DeclareMathOperator{\tamark}{tanh}
\DeclareMathOperator{\sigmoid}{sigmoid}
\DeclareMathOperator{\BCE}{BCE}

% ===========================================================================
\title{DSCSP: Diffusion-Stego with Conditional Sampling Paths for Robust Image Steganography}

\author{
    \IEEEauthorblockN{Authors Omitted for Review}
    \IEEEauthorblockA{}
}

\begin{document}
\maketitle

% ===========================================================================
%  ABSTRACT
% ===========================================================================
\begin{abstract}
This paper introduces DSCSP (Diffusion-Stego with Conditional Sampling Paths), a steganographic framework designed to embed secret messages into images generated by pre-trained diffusion models.
Unlike prior methods that encode information in the initial noise vector or require modification of the generative backbone, DSCSP leaves the diffusion model fully frozen and instead employs a lightweight learned \emph{Controller} to convert payload bits into structured spatial perturbations applied to the output of a deterministic DDIM sampling process.
In doing so, the diffusion model’s sampling trajectory is effectively repurposed as a covert communication channel.
A spatially aligned \emph{Extractor} network based on ResNet-18 then recovers the embedded message from the resulting stego image.

The Controller and Extractor are trained jointly using a composite objective that balances accurate bit reconstruction, perceptual fidelity, and pattern diversity.
To enhance robustness under realistic transmission conditions, we adopt a curriculum-based training strategy that progressively incorporates differentiable approximations of common image distortions, including noise, blur, color adjustments, cropping, and JPEG compression.
Experiments on CelebA-HQ images at $256\times256$ resolution demonstrate that DSCSP embeds 256 bits with near-perfect recovery and minimal perceptual distortion under clean conditions, while retaining resilience against several standard post-processing operations.
We further show that the framework scales from 128 to 512 bits by adjusting the per-cell encoding rate, and we provide an information-theoretic perspective that situates the system’s effective throughput relative to classical channel capacity limits.
\end{abstract}

\begin{IEEEkeywords}
    Steganography, diffusion models, deep learning, image security, robustness
\end{IEEEkeywords}

% ===========================================================================
%  I. INTRODUCTION
% ===========================================================================
\section{Introduction}
\label{sec:intro}

For many years, steganographic methods have focused on modifying existing cover images by subtly altering pixel intensities or transform-domain coefficients, ensuring that the embedded signal remains imperceptible to human observers~\cite{cox2007digital}.
With the emergence of deep generative models, this traditional paradigm has evolved.
In \emph{generative steganography}, the image itself is synthesized with the secret message embedded during the generation process, eliminating reliance on a pre-existing cover image and offering a more inherent form of plausible deniability.


Diffusion models~\cite{ho2020denoising, song2020denoising} now produce images rivaling real photographs in quality, which makes them appealing as the generative backbone for steganographic systems.
Current diffusion-based steganography work can be grouped into two camps.
\emph{Noise-projection methods}, such as Diffusion-Stego~\cite{wei2023diffusionstego}, encode the message into the starting noise vector $\mathbf{z}_T$ and then run the standard reverse process---an elegant training-free idea, but one that provides limited defense against post-processing.
\emph{Model-adaptation methods}~\cite{peng2023stegodiffusion, zhang2024ldstega} modify or extend the diffusion model weights, which is expensive and can degrade the model's original generation quality.

We propose a fundamentally different approach.
\textbf{DSCSP} (Diffusion-Stego with Conditional Sampling Paths) treats the \emph{sampling trajectory} of a frozen pre-trained diffusion model as an information channel.
Rather than modifying the model weights or the initial noise, DSCSP learns a compact \emph{Controller} network that maps payload bits to spatial perturbation patterns, which are applied to the output of the standard DDIM~\cite{song2020denoising} deterministic sampling chain.
A lightweight \emph{Extractor} network recovers the hidden bits from the resulting stego image.
Both networks are trained end-to-end while the diffusion model remains completely frozen, preserving its generation quality and enabling efficient training on consumer-grade hardware.

To ensure robustness against common image transformations encountered during transmission and storage, we introduce a curriculum-based adversarial training strategy that gradually exposes the system to five differentiable attack proxies---Gaussian noise, Gaussian blur, brightness/contrast adjustment, random crop-and-resize, and STE-based differentiable JPEG compression.

Our contributions are as follows:
\begin{itemize}[itemsep=2pt, topsep=2pt]
    \item We introduce \textbf{sampling-path encoding}, a novel steganographic channel for diffusion models that maps bits to learned spatial perturbation patterns applied to frozen DDIM outputs, achieving 256-bit payloads with $0.0000$ BER, $36.7$\,dB PSNR, $0.902$ SSIM, and FID $13.7$ at $256\times256$ resolution.

    \item We propose a \textbf{multi-bit spatial tiling Controller} with orthogonality-regularized patterns that encodes 4 bits per spatial cell using only $\sim$20K trainable parameters, enabling memory-efficient training on a 6\,GB GPU.

    \item We design a \textbf{spatially-aligned Extractor} with a ResNet-18 backbone whose stride-32 feature map naturally aligns with the Controller's $8\times8$ spatial grid, ensuring coherent bit recovery.

    \item We develop a \textbf{differentiable attack curriculum} comprising five attacks with a three-phase training schedule that progressively builds robustness while maintaining reconstruction accuracy.

    \item We provide a \textbf{comprehensive experimental evaluation} including real (non-differentiable) JPEG robustness tests, steganalysis detection experiments with a trained ResNet-18 classifier, capacity scaling from 128 to 512 bits, a baseline comparison with prior methods, and an information-theoretic channel capacity analysis.
\end{itemize}


% ===========================================================================
%  II. RELATED WORK
% ===========================================================================
\section{Related Work}
\label{sec:related}

\subsection{Classical Image Steganography}

Traditional steganography methods embed messages by modifying pixel values~\cite{cox2007digital} or transform-domain coefficients of existing cover images.
LSB (Least Significant Bit) replacement~\cite{chan2004hiding} modifies the least significant bits of pixel values but is easily detectable by statistical steganalysis.
More sophisticated approaches such as HUGO~\cite{pevny2010using}, WOW~\cite{holub2012designing}, and S-UNIWARD~\cite{holub2014universal} minimize a distortion function to reduce detectability, but remain fundamentally limited by the cover image distribution.

\subsection{Deep Learning-Based Steganography}

A second line of work uses neural networks to learn the encoding and decoding functions together.
HiDDeN~\cite{zhu2018hidden}, for instance, pairs an encoder and decoder with a discriminator and a differentiable JPEG proxy so that the hidden message survives compression.
SteganoGAN~\cite{zhang2019steganogan} takes a GAN-based route, folding the message directly into the generation process.
These learning-based approaches significantly increase embedding capacity compared to classical techniques; however, the visual fidelity of the resulting stego images remains inherently constrained by the expressiveness and limitations of the underlying generative backbone.

\subsection{Diffusion-Based Steganography}

Diffusion-Stego~\cite{wei2023diffusionstego} adopts a training-free strategy by embedding the secret message directly into the initial noise vector $\mathbf{z}_T$ using an orthogonal projection mechanism.
The idea is clean, but the encoded signal is spread thinly across the full noise space, leaving it vulnerable to spatial transformations.

StegoDiffusion~\cite{peng2023stegodiffusion} investigates invertible diffusion architectures, exploiting the one-to-one mapping between image and latent space for embedding.
LDStega~\cite{zhang2024ldstega} works within the latent space of latent diffusion models, which gives it practical robustness by operating in a more compact representation.
DERO~\cite{zhao2024dero} tackles robustness head-on by training a message encoder--decoder jointly with the diffusion model, and reports resilience across several classes of image manipulation.

Our work departs from all of these: rather than touching the noise, the latent space, or the model weights, DSCSP attaches a small \emph{learned perturbation controller} to the \emph{output} of a fully frozen diffusion model---keeping the generative backbone untouched and the training footprint small.

\subsection{Differentiable Image Processing for Robustness}

Making a steganographic system robust means training it against image transformations, which in turn requires differentiable stand-ins for non-differentiable operations like JPEG compression.
The Straight-Through Estimator (STE)~\cite{bengio2013estimating} is widely used for this purpose: it lets gradients pass through quantization unchanged, and has been adopted for differentiable JPEG simulation~\cite{shin2017jpeg}.
We build on this idea and combine an STE-based JPEG proxy with differentiable versions of blur, noise, cropping, and color adjustment to form the attack layer used during our robustness training.


% ===========================================================================
%  III. METHODOLOGY
% ===========================================================================
\section{Methodology}
\label{sec:method}

\subsection{System Overview}
\label{sec:overview}

DSCSP consists of three components: (1)~a frozen pre-trained diffusion model $D_\theta$ that generates high-quality base images, (2)~a trainable Controller $C_\phi$ that maps payload bits to a spatial perturbation signal, and (3)~a trainable Extractor $E_\psi$ that recovers the payload from the stego image.

The encoding process works as follows.
Given a payload $\mathbf{b} \in \{0,1\}^{N_b}$ with $N_b = 256$ bits:
\begin{enumerate}[itemsep=1pt, topsep=2pt]
    \item The Controller computes a perturbation: $\bm{\delta} = C_\phi(\mathbf{b}) \in \R^{3 \times H \times W}$.
    \item The diffusion model generates a clean image via DDIM sampling: $\mathbf{x}_\text{clean} = D_\theta(\mathbf{z}_T)$ under \texttt{no\_grad}.
    \item The stego image is formed: $\mathbf{x}_s = \text{clamp}(\mathbf{x}_\text{clean} + \bm{\delta},\, -1,\, 1)$.
\end{enumerate}
For extraction, the Extractor produces logits: $\hat{\mathbf{b}} = \sigmoid(E_\psi(\mathbf{x}_s))$.
Bits are decoded by thresholding at $0.5$.

An important aspect of our design is that we run the full DDIM sampling chain with gradients disabled (\texttt{torch.no\_grad()}), and only apply the perturbation $\bm{\delta}$ \emph{after} the sampling is complete.
This keeps the gradient path simple:
\begin{equation}
    \nabla_\phi \Lcal = \frac{\partial \Lcal}{\partial \hat{\mathbf{b}}} \cdot \frac{\partial \hat{\mathbf{b}}}{\partial \mathbf{x}_s} \cdot \frac{\partial \mathbf{x}_s}{\partial \bm{\delta}} \cdot \frac{\partial \bm{\delta}}{\partial \phi}
    \label{eq:grad_path}
\end{equation}
so gradients never have to propagate through the heavy diffusion UNet, which speeds up training and keeps memory consumption low.

\subsection{Controller Network}
\label{sec:controller}

The role of the Controller is to transform a 256-bit payload $\mathbf{b} \in \{0,1\}^{256}$ into a spatial perturbation map $\bm{\delta} \in \R^{3 \times 256 \times 256}$ using a mechanism we term \emph{multi-bit spatial tiling}.

\subsubsection{Spatial Grid}
To achieve this, we partition the $256 \times 256$ image into an $8 \times 8$ grid, where each cell corresponds to a $32 \times 32$ pixel region, yielding $N_c = 64$ cells in total.
Assigning $K = 4$ bits to each cell allows the grid to encode $N_c \times K = 256$ bits overall.

\subsubsection{Learned Patterns}
Each cell $j \in \{1, \ldots, 64\}$ has $K = 4$ independent learned pattern tensors:
\begin{equation}
    \mathbf{P}_j = \{\mathbf{p}_{j,k}\}_{k=1}^{K}, \quad \mathbf{p}_{j,k} \in \R^{3 \times 32 \times 32}
    \label{eq:patterns}
\end{equation}
All patterns are implemented as a single learnable parameter tensor of shape $(64, 4, 3, 32, 32)$, initialized by sampling from a normal distribution $\Ncal(0, 0.25^2)$.

\subsubsection{Forward Pass}
Given bits $\mathbf{b} = (b_1, \ldots, b_{256})$:
\begin{align}
    &\text{Group:} \quad \mathbf{b}^{(j)} = (b_{4j-3}, \ldots, b_{4j}) \in \{0,1\}^4 \label{eq:group} \\
    &\text{Signs:} \quad s_{j,k} = 2 b^{(j)}_k - 1 \in \{-1, +1\} \label{eq:signs} \\
    &\text{Cell pattern:} \quad \mathbf{c}_j = \sum_{k=1}^{K} s_{j,k} \cdot \mathbf{p}_{j,k} \label{eq:cell} \\
    &\text{Tile:} \quad \bm{\delta}_\text{raw} = \text{tile}(\{\mathbf{c}_j\}_{j=1}^{64}) \in \R^{3 \times 256 \times 256} \label{eq:tile} \\
    &\text{Constrain:} \quad \bm{\delta} = \tanh(\bm{\delta}_\text{raw}) \cdot \alpha_\text{max} \label{eq:constrain}
\end{align}
where $\alpha_{\text{max}} = 0.3$ defines the maximum allowable perturbation magnitude.
The $\tanh$ nonlinearity constrains the perturbation to the range $[-\alpha_{\text{max}}, \alpha_{\text{max}}]$, thereby limiting the maximum change applied to any pixel.
Because each cell combines $K = 4$ learned patterns, the resulting effective signal-to-noise ratio per bit is $\frac{1}{K-1} = \frac{1}{3}$, which remains sufficient for the Extractor to disentangle and recover the individual bit contributions reliably.

The Controller has approximately \textbf{20K trainable parameters} ($64 \times 4 \times 3 \times 32 \times 32 = 786{,}432$ floating-point values in the pattern tensor).

\subsection{Extractor Network}
\label{sec:extractor}

Given a stego image $\mathbf{x}_s \in \R^{3 \times 256 \times 256}$, the goal of the Extractor is to reconstruct the full 256-bit payload embedded within it.

\subsubsection{Architecture}
We adopt a ResNet-18~\cite{he2016deep} backbone initialized from scratch, removing the final average pooling and fully connected layers.
For an input of size $256 \times 256$, the network’s cumulative stride of 32 yields feature maps of shape $(B, 512, 8, 8)$, which align naturally with the Controller’s $8\times8$ spatial tiling structure.

A $1\times1$ convolutional head reduces the 512 channels to $K = 4$ channels:
\begin{equation}
    \mathbf{f} = \text{Conv}_{1\times1}^{512 \to 4}(\text{ResNet18}(\mathbf{x}_s)) \in \R^{B \times 4 \times 8 \times 8}
    \label{eq:extractor_head}
\end{equation}

\subsubsection{Bit Ordering Alignment}
To guarantee that the recovered bits follow the same ordering used by the Controller during encoding, the feature map is permuted prior to flattening:
\begin{equation}
    \hat{\mathbf{b}}_\text{logits} = \text{reshape}\big(\text{permute}(\mathbf{f},\, (0,2,3,1)),\, (B, 256)\big)
    \label{eq:permute}
\end{equation}
This permutation reshuffles the tensor from $(B, K, 8, 8)$ to $(B, 8, 8, K)$ before flattening. The result is that the $i$-th output logit maps to cell $\lfloor i/K \rfloor$ at within-cell position $i \bmod K$, which lines up exactly with how the Controller organizes its \texttt{view(B, 64, 4)} bit layout.

The Extractor has approximately \textbf{11.2M trainable parameters}.

\subsection{Frozen Diffusion Backbone}
\label{sec:diffusion}

For our generative backbone, we use \texttt{google/ddpm-celebahq-256}, a pre-trained DDPM~\cite{ho2020denoising} available through the HuggingFace \texttt{diffusers} library, which produces $256\times256$ face images.

We load the UNet in FP16 and freeze every parameter ($\texttt{requires\_grad\_ = False}$).
Image generation follows the deterministic DDIM~\cite{song2020denoising} schedule with $T = 20$ steps.
Given initial noise $\mathbf{z}_T \sim \Ncal(\mathbf{0}, \mathbf{I})$, the DDIM update at step $i$ with timestep $t$ is:
\begin{align}
    \hat{\bm{\epsilon}} &= D_\theta(\mathbf{z}_t, t) \label{eq:noise_pred} \\
    \hat{\mathbf{x}}_0 &= \frac{\mathbf{z}_t - \sqrt{1 - \bar{\alpha}_t}\, \hat{\bm{\epsilon}}}{\sqrt{\bar{\alpha}_t}} \label{eq:x0_pred} \\
    \mathbf{z}_{t-1} &= \sqrt{\bar{\alpha}_{t-1}}\, \hat{\mathbf{x}}_0 + \sqrt{1 - \bar{\alpha}_{t-1}}\, \hat{\bm{\epsilon}} \label{eq:ddim_step}
\end{align}
where $\bar{\alpha}_t$ denotes the cumulative product of the noise schedule coefficients.

\subsubsection{Memory Optimization}
Since we target a 6\,GB GPU, we pre-generate a pool of 128 clean images and keep them in GPU memory, then move the UNet back to CPU---reclaiming roughly 440\,MB of VRAM.
Every 2{,}000 training steps the pool is regenerated so the model does not memorize a fixed set of covers.

\subsection{Loss Functions}
\label{sec:losses}

The total training objective is:
\begin{equation}
    \Lcal = \Lcal_\text{msg} + \lambda_p \Lcal_\text{perc} + \lambda_c \Lcal_\text{ctrl} + \lambda_o \Lcal_\text{ortho}
    \label{eq:total_loss}
\end{equation}

\subsubsection{Message Reconstruction Loss}
Binary cross-entropy with logits between predicted and ground-truth bits:
\begin{equation}
    \Lcal_\text{msg} = \BCE(\hat{\mathbf{b}}_\text{logits},\, \mathbf{b})
    \label{eq:loss_msg}
\end{equation}

\subsubsection{Perceptual Fidelity Loss}
LPIPS~\cite{zhang2018unreasonable} distance using a frozen VGG backbone:
\begin{equation}
    \Lcal_\text{perc} = \text{LPIPS}_\text{VGG}(\mathbf{x}_s,\, \mathbf{x}_\text{clean})
    \label{eq:loss_perc}
\end{equation}

\subsubsection{Controller Regularization}
Mean squared perturbation magnitude:
\begin{equation}
    \Lcal_\text{ctrl} = \frac{1}{3HW} \|\bm{\delta}\|_F^2
    \label{eq:loss_ctrl}
\end{equation}

\subsubsection{Orthogonality Regularization}
To prevent the $K$ learned patterns within each cell from collapsing into dependent directions, we penalize deviation of the Gram matrix from the identity:
\begin{equation}
    \Lcal_\text{ortho} = \frac{1}{N_c} \sum_{j=1}^{N_c} \left\| \mathbf{G}_j - \mathbf{I}_K \right\|_F^2
    \label{eq:loss_ortho}
\end{equation}
where $\mathbf{G}_j \in \R^{K \times K}$ is the Gram matrix of the $K$ flattened patterns in cell $j$:
\begin{equation}
    (\mathbf{G}_j)_{k,l} = \langle \text{vec}(\mathbf{p}_{j,k}),\, \text{vec}(\mathbf{p}_{j,l}) \rangle
\end{equation}

Loss weights are set to $\lambda_p = 0.1$, $\lambda_c = 0.0$ (disabled), and $\lambda_o = 0.001$.

\subsection{Differentiable Attack Layer}
\label{sec:attacks}

To train for robustness against common image transformations, we implement five differentiable attacks that preserve gradient flow back to the Controller.

\subsubsection{Gaussian Noise}
Additive zero-mean Gaussian noise with $\sigma \sim \Ucal(0.01, 0.05)$:
\begin{equation}
    \mathbf{x}' = \text{clamp}(\mathbf{x} + \bm{\epsilon},\, -1, 1), \quad \bm{\epsilon} \sim \Ncal(\mathbf{0}, \sigma^2 \mathbf{I})
\end{equation}

\subsubsection{Gaussian Blur}
Depthwise convolution with a $5\times5$ Gaussian kernel ($\sigma \sim \Ucal(0.5, 2.0)$):
\begin{equation}
    \mathbf{x}' = \mathbf{x} \ast \mathbf{K}_\sigma, \quad K_\sigma(i,j) \propto \exp\!\left(-\frac{i^2 + j^2}{2\sigma^2}\right)
\end{equation}

\subsubsection{Brightness/Contrast}
Affine intensity transformation with $c \sim \Ucal(0.8, 1.2)$, $\beta \sim \Ucal(-0.1, 0.1)$:
\begin{equation}
    \mathbf{x}' = \text{clamp}(c \cdot \mathbf{x} + \beta,\, -1, 1)
\end{equation}

\subsubsection{Crop and Resize}
Random spatial crop with ratio $r \sim \Ucal(0.7, 0.95)$ followed by bilinear interpolation to the original resolution:
\begin{equation}
    \mathbf{x}' = \text{resize}_\text{bilinear}(\text{crop}_r(\mathbf{x}),\, H \times W)
\end{equation}

\subsubsection{Differentiable JPEG (STE)}
A full JPEG simulation pipeline: RGB $\to$ YCbCr color conversion, $8\times8$ block DCT via matrix multiplication, quantization using standard luminance and chrominance tables scaled by quality factor $Q \sim \{50, \ldots, 95\}$, block IDCT, and YCbCr $\to$ RGB conversion.
The quantization step uses the Straight-Through Estimator~\cite{bengio2013estimating}:
\begin{equation}
    \hat{d} = d + \left(\text{round}\!\left(\frac{d}{q}\right) \cdot q - d\right)_\text{detach}
    \label{eq:ste}
\end{equation}
where $d$ denotes DCT coefficients and $q$ the quantization step size.
In the forward pass, $\hat{d}$ equals the quantized value; in the backward pass, $\nabla_d \hat{d} = 1$ (identity gradient).
The quality factor scaling follows the standard JPEG specification:
\begin{equation}
    s = \begin{cases} 5000 / Q & \text{if } Q < 50 \\ 200 - 2Q & \text{if } Q \geq 50 \end{cases}
    \label{eq:jpeg_scale}
\end{equation}

\subsubsection{Attack Wrapper}
During training, the \texttt{AttackLayer} module randomly selects one of the five attacks per batch with probability $p_\text{atk}$ (determined by the curriculum schedule), or passes the image through unchanged with probability $1 - p_\text{atk}$.

\subsection{Training Procedure}
\label{sec:training}

\subsubsection{Image Pool Strategy}
Rather than running the diffusion model at every training step, we pre-generate a pool of $M = 128$ clean images using DDIM sampling and cache them in GPU memory.
At each training step, a random batch of $B = 8$ images is sampled from the pool.
The pool is regenerated every 2{,}000 steps to prevent overfitting to a fixed set of images.
After pool generation, the UNet is offloaded to CPU to free GPU memory.

\subsubsection{Robustness Curriculum}
Training is organized into three phases with a curriculum for the attack probability $p_\text{atk}$:
\begin{equation}
    p_\text{atk}(t) = \begin{cases}
        0 & \text{if } t < 3{,}000 \\
        0.7 \cdot \frac{t - 3{,}000}{7{,}000} & \text{if } 3{,}000 \leq t < 10{,}000 \\
        0.7 & \text{if } t \geq 10{,}000
    \end{cases}
    \label{eq:curriculum}
\end{equation}

\begin{itemize}[itemsep=1pt, topsep=2pt]
    \item \textbf{Phase~1} (steps $0$--$3{,}000$): No attacks are introduced. During this stage, the Controller and Extractor concentrate solely on learning the core encoding and decoding mapping.
    \item \textbf{Phase~2} (steps $3{,}000$--$10{,}000$): The attack probability increases linearly from $0$ to $0.7$, enabling the networks to progressively adapt to distortions in a controlled manner.
    \item \textbf{Phase~3} (steps $10{,}000$--$30{,}000$): The attack probability is fixed at $0.7$, so most training batches undergo a randomly selected transformation, reinforcing robustness.
\end{itemize}

\subsubsection{Optimization}
Training is performed using Adam~\cite{kingma2014adam} with $\beta_1 = 0.9$ and $\beta_2 = 0.999$.
The Controller and Extractor are assigned separate parameter groups: the Controller is trained without weight decay, while the Extractor uses a decay factor of $10^{-4}$.
Both groups are initialized with a learning rate of $5 \times 10^{-4}$, which is annealed to $10^{-5}$ over 30{,}000 steps using a cosine schedule~\cite{loshchilov2017sgdr}.
Gradients are clipped to a maximum norm of $5.0$ to ensure stable optimization.

Algorithms~\ref{alg:encoding} and~\ref{alg:extraction} provide a concise summary of the encoding and extraction procedures.

\begin{algorithm}[t]
\caption{DSCSP Encoding}
\label{alg:encoding}
\begin{algorithmic}[1]
\REQUIRE Payload bits $\mathbf{b} \in \{0,1\}^{256}$, frozen diffusion model $D_\theta$, and trained Controller $C_\phi$
\ENSURE Stego image $\mathbf{x}_s \in [-1, 1]^{3 \times 256 \times 256}$
\STATE $\bm{\delta} \leftarrow C_\phi(\mathbf{b})$ \COMMENT{Compute perturbation}
\STATE $\mathbf{z}_T \sim \Ncal(\mathbf{0}, \mathbf{I})$ \COMMENT{Sample initial noise}
\STATE $\mathbf{z} \leftarrow \mathbf{z}_T$
\FOR{$i = 1$ \TO $T$}
    \STATE $\hat{\bm{\epsilon}} \leftarrow D_\theta(\mathbf{z}, t_i)$ \COMMENT{Predict noise (no grad)}
    \STATE $\hat{\mathbf{x}}_0 \leftarrow (\mathbf{z} - \sqrt{1 - \bar{\alpha}_{t_i}}\, \hat{\bm{\epsilon}}) / \sqrt{\bar{\alpha}_{t_i}}$
    \STATE $\mathbf{z} \leftarrow \sqrt{\bar{\alpha}_{t_{i-1}}}\, \hat{\mathbf{x}}_0 + \sqrt{1 - \bar{\alpha}_{t_{i-1}}}\, \hat{\bm{\epsilon}}$
\ENDFOR
\STATE $\mathbf{x}_\text{clean} \leftarrow \text{clamp}(\mathbf{z}, -1, 1)$
\STATE $\mathbf{x}_s \leftarrow \text{clamp}(\mathbf{x}_\text{clean} + \bm{\delta}, -1, 1)$
\RETURN $\mathbf{x}_s$
\end{algorithmic}
\end{algorithm}

\begin{algorithm}[t]
\caption{DSCSP Extraction}
\label{alg:extraction}
\begin{algorithmic}[1]
\REQUIRE Stego image $\mathbf{x}_s \in [-1, 1]^{3 \times 256 \times 256}$, trained Extractor $E_\psi$
\ENSURE Recovered bits $\hat{\mathbf{b}} \in \{0,1\}^{256}$
\STATE $\mathbf{f} \leftarrow \text{ResNet18}(\mathbf{x}_s)$ \COMMENT{Feature extraction: $(B, 512, 8, 8)$}
\STATE $\mathbf{h} \leftarrow \text{Conv}_{1\times1}^{512 \to 4}(\mathbf{f})$ \COMMENT{Spatial head: $(B, 4, 8, 8)$}
\STATE $\mathbf{h} \leftarrow \text{permute}(\mathbf{h}, (0,2,3,1))$ \COMMENT{Reorder: $(B, 8, 8, 4)$}
\STATE $\hat{\mathbf{b}}_\text{logits} \leftarrow \text{reshape}(\mathbf{h}, (B, 256))$
\STATE $\hat{\mathbf{b}} \leftarrow \mathbb{1}[\sigmoid(\hat{\mathbf{b}}_\text{logits}) > 0.5]$ \COMMENT{Threshold}
\RETURN $\hat{\mathbf{b}}$
\end{algorithmic}
\end{algorithm}


% ===========================================================================
%  IV. EXPERIMENTS
% ===========================================================================
\section{Experiments}
\label{sec:experiments}

\subsection{Experimental Setup}
\label{sec:setup}

\subsubsection{Hardware}
All experiments are performed on an NVIDIA RTX 3060 Laptop GPU equipped with 6\,GB of VRAM, using PyTorch 2.0+ with CUDA acceleration.

\subsubsection{Diffusion Model}
We use the pre-trained \texttt{google/ddpm-celebahq-256} DDPM from HuggingFace, generating $256\times256$ face images via 20-step DDIM sampling.
The UNet has approximately 114M parameters, all frozen in FP16.

\subsubsection{Configuration}
Table~\ref{tab:hyperparams} summarizes all hyperparameters.

\begin{table}[t]
\centering
\caption{Training hyperparameters.}
\label{tab:hyperparams}
\small
\begin{tabular}{@{}llr@{}}
\toprule
\textbf{Category} & \textbf{Parameter} & \textbf{Value} \\
\midrule
\multirow{4}{*}{Model}
    & Payload capacity & 256 bits \\
    & Bits per cell & 4 \\
    & Grid size & $8 \times 8$ \\
    & Max perturbation $\alpha_\text{max}$ & 0.3 \\
\midrule
\multirow{6}{*}{Training}
    & Total steps & 30{,}000 \\
    & Batch size & 8 \\
    & Learning rate & $5 \times 10^{-4}$ \\
    & LR schedule & Cosine $\to 10^{-5}$ \\
    & Gradient clip & 5.0 \\
    & Weight decay (Extractor) & $10^{-4}$ \\
\midrule
\multirow{3}{*}{Loss}
    & $\lambda_p$ (LPIPS) & 0.1 \\
    & $\lambda_c$ (Controller L2) & 0.0 \\
    & $\lambda_o$ (Orthogonality) & 0.001 \\
\midrule
\multirow{2}{*}{Pool}
    & Pool size & 128 \\
    & Refresh interval & 2{,}000 steps \\
\midrule
\multirow{3}{*}{Robustness}
    & Curriculum start & Step 3{,}000 \\
    & Ramp duration & 7{,}000 steps \\
    & Max attack prob. & 0.7 \\
\midrule
\multirow{5}{*}{Attack ranges}
    & JPEG quality & $\{50, \ldots, 95\}$ \\
    & Blur $\sigma$ & $\Ucal(0.5, 2.0)$ \\
    & Noise $\sigma$ & $\Ucal(0.01, 0.05)$ \\
    & Crop ratio & $\Ucal(0.7, 0.95)$ \\
    & Brightness/contrast & $\pm 0.1$ / $\times[0.8, 1.2]$ \\
\bottomrule
\end{tabular}
\end{table}

\subsubsection{Evaluation Metrics}
We report the following metrics:
\begin{itemize}[itemsep=1pt, topsep=2pt]
    \item \textbf{Bit Error Rate (BER):} The proportion of bits that are incorrectly decoded. A value of $\text{BER} = 0$ corresponds to perfect recovery, whereas $\text{BER} = 0.5$ reflects performance equivalent to random guessing.
    \item \textbf{Peak Signal-to-Noise Ratio (PSNR):} Measures pixel-level fidelity between clean and stego images. For images in $[-1, 1]$:
    \begin{equation}
        \text{PSNR} = 10 \log_{10}\!\left(\frac{4}{\text{MSE}(\mathbf{x}_s, \mathbf{x}_\text{clean})}\right)
    \end{equation}
    Higher values indicate less visible distortion.
    \item \textbf{Structural Similarity (SSIM)}~\cite{wang2004image}: Measures perceptual structural similarity. SSIM $\in [0, 1]$, where 1 indicates identical images.
    \item \textbf{LPIPS}~\cite{zhang2018unreasonable}: Learned Perceptual Image Patch Similarity computed using a VGG-16 backbone. Lower values indicate higher perceptual similarity.
    \item \textbf{Fr\'{e}chet Inception Distance (FID)}~\cite{heusel2017gans}: Measures the distributional distance between clean and stego image populations using Inception-v3 features. Lower FID indicates statistical indistinguishability.
\end{itemize}

\subsection{Training Dynamics}
\label{sec:dynamics}

Table~\ref{tab:training_dynamics} reports key metrics at milestone steps throughout training.

\begin{table}[t]
\centering
\caption{Training dynamics at key milestones. BER is measured on clean (non-attacked) batches. The ``Atk'' column shows the attack probability at each step.}
\label{tab:training_dynamics}
\small
\begin{tabular}{@{}rcccc@{}}
\toprule
\textbf{Step} & \textbf{BER} & \textbf{PSNR (dB)} & \textbf{Atk} & \textbf{$\Lcal_\text{ortho}$} \\
\midrule
0 (init)     & 0.4688 & 24.1 & 0.00 & 8{,}728 \\
440          & 0.0000 & 30.0 & 0.00 & --- \\
1{,}000      & 0.0000 & 30.7 & 0.00 & --- \\
2{,}000      & 0.0000 & 34.2 & 0.00 & $\sim$10 \\
3{,}000      & 0.0000 & 36.8 & 0.00 & $\sim$10 \\
4{,}000      & 0.0005 & 37.6 & 0.10 & --- \\
5{,}000      & 0.0000 & 36.8 & 0.20 & --- \\
7{,}000      & 0.0000 & 35.0 & 0.40 & --- \\
10{,}000     & 0.0000 & 35.2 & 0.70 & $\sim$25 \\
15{,}000     & 0.0000 & 35.9 & 0.70 & $\sim$18 \\
20{,}000     & 0.0000 & 36.3 & 0.70 & $\sim$14 \\
25{,}000     & 0.0000 & 36.5 & 0.70 & $\sim$12 \\
\textbf{30{,}000}     & \textbf{0.0000} & \textbf{36.7} & \textbf{0.70} & \textbf{11.4} \\
\bottomrule
\end{tabular}
\end{table}

\textit{Phase~1: Clean Convergence (Steps 0--3{,}000).}
Training converges rapidly---the BER falls to zero by step 440, indicating that the Controller--Extractor pair learns the fundamental encode--decode mapping within minutes.
During this interval, PSNR increases from 24.1\,dB to 36.8\,dB as the Controller progressively reduces perturbation magnitude while maintaining reliable message recovery.
Meanwhile, the orthogonality loss $\Lcal_\text{ortho}$ decreases sharply from 8{,}728 to approximately 10, confirming that the per-cell patterns evolve toward near-orthogonal configurations.

\textit{Phase~2: Robustness Ramp (Steps 3{,}000--10{,}000).}
As attacks are gradually introduced, a clear quality--robustness tradeoff emerges: PSNR temporarily declines to around 35.0\,dB because the Controller must strengthen its perturbations for resilience rather than subtlety alone.
Notably, the clean BER remains at zero throughout this stage, demonstrating that the underlying encoding capacity remains intact despite the added distortions.

\textit{Phase~3: Full Robustness (Steps 10{,}000--30{,}000).}
With the attack probability fixed at 70\%, the system enters a stable regime.
As the cosine learning-rate schedule progresses, PSNR steadily recovers to 36.7\,dB while the Controller refines its perturbation patterns.
The orthogonality loss continues to decline, reaching 11.4 by the end of training.
The final model achieves perfect bit recovery (BER = 0.0000) on clean images while sustaining a PSNR of 36.7\,dB.

\subsection{Main Results}
\label{sec:results}

\subsubsection{Clean Performance}
Table~\ref{tab:main_results} reports the final evaluation results on 500 newly generated test images.
Across all reported metrics, DSCSP achieves perfect bit recovery while maintaining low perceptual distortion.

\begin{table}[t]
\centering
\caption{Comprehensive evaluation on 500 random samples (clean, no attacks). All metrics computed using real (non-differentiable) implementations.}
\label{tab:main_results}
\begin{tabular}{@{}lcc@{}}
\toprule
\textbf{Metric} & \textbf{Target} & \textbf{Achieved} \\
\midrule
Bit Error Rate (BER)        & $< 0.05$ & $\mathbf{0.0000}$ \\
PSNR (dB)                   & $> 25$   & $\mathbf{36.72}$ \\
SSIM                        & $> 0.85$ & $\mathbf{0.902}$ \\
LPIPS (VGG)                 & $< 0.10$ & $\mathbf{0.032}$ \\
FID                         & ---      & $\mathbf{13.7}$ \\
\midrule
Payload capacity             & ---      & 256 bits \\
Bits per pixel (bpp)         & ---      & $3.91 \times 10^{-3}$ \\
Controller parameters        & ---      & $\sim$20K \\
Extractor parameters         & ---      & $\sim$11.2M \\
Training time                & ---      & $\sim$100 min \\
\bottomrule
\end{tabular}
\end{table}

\subsubsection{Robustness Against Real Attacks}
To assess robustness beyond the differentiable proxies used during training, we evaluate the model using real (non-differentiable) implementations of each attack type with standard libraries (e.g., Pillow and NumPy).
Table~\ref{tab:robustness} presents the resulting BER under each attack across 200 test pairs.

\begin{table}[t]
\centering
\caption{Robustness evaluation under real (non-differentiable) attacks on 200 test samples. BER values below $0.01$ are considered indicative of reliable message recovery.}
\label{tab:robustness}
\small
\begin{tabular}{@{}llcc@{}}
\toprule
\textbf{Attack} & \textbf{Parameter} & \textbf{BER} & \textbf{Status} \\
\midrule
None (clean) & --- & 0.0000 & \checkmark \\
\midrule
\multirow{6}{*}{JPEG}
    & $Q = 95$ & 0.0007 & \checkmark \\
    & $Q = 90$ & 0.0075 & \checkmark \\
    & $Q = 80$ & 0.1140 & $\sim$ \\
    & $Q = 70$ & 0.2748 & $\times$ \\
    & $Q = 60$ & 0.3527 & $\times$ \\
    & $Q = 50$ & 0.3639 & $\times$ \\
\midrule
\multirow{4}{*}{Gaussian Noise}
    & $\sigma = 0.01$ & 0.0000 & \checkmark \\
    & $\sigma = 0.02$ & 0.0000 & \checkmark \\
    & $\sigma = 0.03$ & 0.0000 & \checkmark \\
    & $\sigma = 0.05$ & 0.0000 & \checkmark \\
\midrule
\multirow{4}{*}{Gaussian Blur}
    & $r = 0.5$ & 0.0000 & \checkmark \\
    & $r = 1.0$ & 0.0012 & \checkmark \\
    & $r = 1.5$ & 0.0081 & \checkmark \\
    & $r = 2.0$ & 0.0275 & $\sim$ \\
\midrule
\multirow{5}{*}{Bright./Contr.}
    & $b{=}1.1, c{=}1.0$ & 0.0013 & \checkmark \\
    & $b{=}0.9, c{=}1.0$ & 0.0000 & \checkmark \\
    & $b{=}1.0, c{=}1.2$ & 0.0004 & \checkmark \\
    & $b{=}1.0, c{=}0.8$ & 0.0000 & \checkmark \\
    & $b{=}1.1, c{=}1.2$ & 0.0020 & \checkmark \\
\midrule
\multirow{4}{*}{Crop + Resize}
    & ratio $= 0.95$ & 0.4773 & $\times$ \\
    & ratio $= 0.90$ & 0.4649 & $\times$ \\
    & ratio $= 0.80$ & 0.4806 & $\times$ \\
    & ratio $= 0.70$ & 0.4959 & $\times$ \\
\bottomrule
\end{tabular}
\end{table}

Gaussian noise introduces no bit errors across all tested $\sigma$ values, and brightness/contrast adjustments consistently remain below BER $0.003$ under every configuration evaluated.
Gaussian blur is tolerated effectively up to a kernel radius of $1.5$ (BER $< 0.01$), although stronger blurring begins to impair reliable recovery.

\subsubsection{STE-to-Real JPEG Gap}
An important observation is the gap between performance under differentiable STE-based JPEG during training (BER $\sim 0.002$) and real JPEG compression at test time.
For high quality factors ($Q \geq 90$), real JPEG exhibits strong robustness (BER $< 0.01$), indicating that the STE proxy transfers well in this regime.
However, at lower quality settings ($Q \leq 70$), BER increases substantially to $0.27$--$0.36$, revealing limitations of the proxy under aggressive compression.
This discrepancy arises because practical JPEG implementations introduce effects such as chroma subsampling, block-boundary artifacts, and codec-specific rounding behavior that are not fully captured by the differentiable approximation.

\subsubsection{Crop Vulnerability}
Crop-and-resize represents the most severe failure case (BER $\sim 0.47$--$0.50$, effectively random performance), as cropping disrupts the spatial correspondence between the Controller’s grid structure and the Extractor’s feature alignment.

\subsubsection{Model Efficiency}
In terms of parameter overhead, the additional cost is modest.
The Controller contains only $\sim$20K trainable parameters—approximately four orders of magnitude fewer than the frozen 114M-parameter UNet backbone.
Even when including the Extractor, the total trainable parameters ($\sim$11.2M) remain roughly $10\times$ smaller than the backbone.
End-to-end training completes in approximately 100 minutes on a single laptop GPU with 6\,GB VRAM, making the approach accessible without large-scale computational resources.


\subsection{Baseline Comparison}
\label{sec:baselines}

Table~\ref{tab:baselines} places DSCSP alongside several prior methods, using the numbers reported in their respective papers.
We note that a strict apples-to-apples comparison is difficult since each method was developed on different datasets, resolutions, and evaluation setups, so the table should be read with that caveat in mind.

\begin{table}[t]
\centering
\caption{Comparison with prior methods. Published numbers from original papers. ``---'' indicates not reported. Robustness column indicates whether the method was trained/evaluated with image attacks.}
\label{tab:baselines}
\small
\begin{tabular}{@{}lccccc@{}}
\toprule
\textbf{Method} & \textbf{Bits} & \textbf{BER} & \textbf{PSNR} & \textbf{SSIM} & \textbf{Rob.} \\
\midrule
HiDDeN~\cite{zhu2018hidden}         & 30   & $\sim$0.02 & 33.0  & ---   & \checkmark \\
SteganoGAN~\cite{zhang2019steganogan} & $\sim$4 bpp & ---  & 28--32 & ---   & $\times$ \\
Diff.-Stego~\cite{wei2023diffusionstego} & 256  & $\sim$0.0 & ---   & ---   & $\times$ \\
LDStega~\cite{zhang2024ldstega}       & 256  & $< 0.01$ & 30.5  & 0.93  & \checkmark \\
\midrule
\textbf{DSCSP (ours)}               & \textbf{256} & \textbf{0.000} & \textbf{36.7} & \textbf{0.902} & \checkmark \\
\bottomrule
\end{tabular}
\end{table}

Among methods that handle at least 256 bits, DSCSP posts the highest PSNR.
HiDDeN does provide robustness, but only for 30 bits---an order of magnitude fewer.
Diffusion-Stego can match our bit count without any training, yet it has no mechanism to survive post-processing.
LDStega, which embeds in the latent space of latent diffusion models, reports a slightly higher SSIM but at a lower PSNR.
What sets DSCSP apart is that it brings together high capacity, high fidelity, and robustness training in a package that runs on a consumer GPU.


\subsection{Steganalysis Detection}
\label{sec:steganalysis}

To evaluate the statistical detectability of DSCSP stego images, we train a binary classifier to distinguish clean from stego images.

\subsubsection{Setup}
We generate 500 clean/stego image pairs using the trained model and split them into 800 training samples and 200 test samples (balanced classes).
A ResNet-18~\cite{he2016deep} classifier is trained for 30 epochs with Adam ($\text{lr} = 10^{-4}$) and cross-entropy loss.

\subsubsection{Results}
Table~\ref{tab:steganalysis} reports detection metrics on the held-out test set.

\begin{table}[t]
\centering
\caption{Steganalysis detection results. A ResNet-18 classifier trained to distinguish clean from stego images.}
\label{tab:steganalysis}
\begin{tabular}{@{}lc@{}}
\toprule
\textbf{Metric} & \textbf{Value} \\
\midrule
ROC AUC    & 1.0000 \\
Accuracy   & 1.0000 \\
Precision  & 1.0000 \\
Recall     & 1.0000 \\
\bottomrule
\end{tabular}
\end{table}

The detector attains perfect classification performance (AUC = 1.0), indicating that DSCSP stego images are \emph{readily detectable} by a learned classifier.
This outcome is unsurprising given the structured nature of the Controller’s spatial perturbations: the $8 \times 8$ tiled grid introduces consistent artifacts that a sufficiently expressive model can reliably exploit.
Such detectability is a well-documented limitation of many high-capacity steganographic systems~\cite{boroumand2019deep} and reflects the inherent tradeoff between payload capacity and statistical undetectability.
We outline possible mitigation strategies in Section~\ref{sec:discussion}.


\subsection{Capacity Scaling}
\label{sec:scaling}

We investigate how DSCSP's performance scales with payload capacity by varying the bits per cell parameter while keeping the $8 \times 8$ spatial grid fixed.
Three configurations are evaluated: 128 bits ($K=2$ bits/cell), 256 bits ($K=4$, baseline), and 512 bits ($K=8$).
The 128-bit and 512-bit variants are trained for 15{,}000 steps on clean images without robustness attacks.
The 256-bit model uses the fully robustness-trained model (30K steps with attack curriculum), so its PSNR is lower due to the stronger perturbation patterns required for attack resilience.

\begin{table}[t]
\centering
\caption{Capacity scaling: quality--capacity tradeoff at different payload sizes. 128-bit and 512-bit models trained for 15K steps (clean, no attacks). 256-bit model is the full 30K robustness-trained model.}
\label{tab:scaling}
\small
\begin{tabular}{@{}cccccc@{}}
\toprule
\textbf{Bits} & \textbf{$K$} & \textbf{BER} & \textbf{PSNR (dB)} & \textbf{SSIM} & \textbf{LPIPS} \\
\midrule
128  & 2 & 0.0000 & 48.34 & 0.992 & 0.0004 \\
256  & 4 & 0.0000 & 36.72 & 0.902 & 0.032 \\
512  & 8 & 0.0000 & 41.60 & 0.962 & 0.0011 \\
\bottomrule
\end{tabular}
\end{table}

All three capacity variants achieve perfect bit recovery (BER = 0.0000), demonstrating that the spatially-aligned Controller--Extractor architecture reliably supports encoding rates from 2 to 8 bits per cell.
The 128-bit and 512-bit clean-only results show the expected capacity--quality tradeoff: PSNR decreases from 48.34\,dB to 41.60\,dB as the payload quadruples.
The per-cell signal-to-noise ratio decreases as $\text{SNR} = 1/(K-1)$: from $1.0$ at $K=2$ to $0.143$ at $K=8$, yet the Extractor's non-linear decoding maintains zero-error recovery.


\subsection{Theoretical Channel Capacity}
\label{sec:capacity}

We provide an information-theoretic analysis of the DSCSP encoding scheme, framing it as a structured additive channel.

\subsubsection{Channel Model}
Each cell in the $8 \times 8$ grid can be modeled as an independent additive channel.
The Controller maps $K$ bits to a cell perturbation $\mathbf{c}_j = \sum_{k=1}^{K} s_{j,k} \cdot \mathbf{p}_{j,k}$ (Eq.~\ref{eq:cell}), where $s_{j,k} \in \{-1, +1\}$ encodes the $k$-th bit.
If the $K$ patterns $\{\mathbf{p}_{j,k}\}$ are mutually orthogonal, the Extractor can project the received signal onto each pattern independently using a matched filter.

\subsubsection{Per-Cell SNR}
For a given bit $k$ in cell $j$, the desired signal component is $s_{j,k} \cdot \mathbf{p}_{j,k}$, while the interference from other bits is $\sum_{l \neq k} s_{j,l} \cdot \mathbf{p}_{j,l}$.
Under the orthogonality assumption ($\langle \mathbf{p}_{j,k}, \mathbf{p}_{j,l} \rangle = 0$ for $k \neq l$), the matched-filter output for bit $k$ is:
\begin{equation}
    y_k = \langle \mathbf{c}_j, \mathbf{p}_{j,k} \rangle = s_{j,k} \|\mathbf{p}_{j,k}\|^2
\end{equation}
which is free of inter-bit interference.
In the non-orthogonal case, the cross-correlation terms act as noise.
With unit-norm patterns, the signal power per bit is $1$ and the total interference power from $K-1$ other bits is at most $K-1$, yielding:
\begin{equation}
    \text{SNR}_\text{per-cell} = \frac{1}{K - 1}
\end{equation}

\subsubsection{Shannon Capacity Bound}
For an additive white Gaussian noise (AWGN) channel, the Shannon capacity per channel use is $C = \frac{1}{2}\log_2(1 + \text{SNR})$.
Across $N_c = 64$ cells, each encoding $K$ bits, the maximum reliable throughput is bounded by:
\begin{equation}
    C_\text{total} = N_c \cdot K \cdot \frac{1}{2} \log_2\!\left(1 + \frac{1}{K-1}\right) \;\text{bits}
    \label{eq:shannon}
\end{equation}

Table~\ref{tab:capacity_theory} compares the achieved capacity to Shannon bounds.

\begin{table}[t]
\centering
\caption{Theoretical channel capacity analysis. $C_\text{Shannon}$ computed from Eq.~\ref{eq:shannon}. Efficiency = Achieved bits / $C_\text{Shannon}$.}
\label{tab:capacity_theory}
\small
\begin{tabular}{@{}ccccc@{}}
\toprule
\textbf{$K$} & \textbf{SNR} & \textbf{Payload} & \textbf{$C_\text{Shannon}$} & \textbf{Efficiency} \\
\midrule
2 & 1.000 & 128 & 64.0 & 200\% \\
4 & 0.333 & 256 & 53.7 & 477\% \\
8 & 0.143 & 512 & 61.7 & 830\% \\
\bottomrule
\end{tabular}
\end{table}

\subsubsection{Exceeding the AWGN Bound}
The achieved throughput substantially exceeds the per-cell AWGN Shannon bound.
This is not a violation of Shannon's theorem, but rather reflects the fact that the AWGN model is overly conservative for the DSCSP system.
Three key factors explain the gap:
\begin{enumerate}[itemsep=1pt, topsep=2pt]
    \item \textbf{Learned features, not matched filters.} The Extractor uses a deep ResNet-18 with 11.2M parameters---far more powerful than a linear matched filter. It learns non-linear decision boundaries that can exploit the full spatial structure of patterns.
    \item \textbf{Joint optimization.} The Controller and Extractor are trained jointly, allowing patterns to evolve toward configurations that are maximally separable by the Extractor, rather than constrained to orthogonality.
    \item \textbf{High-dimensional patterns.} Each pattern $\mathbf{p}_{j,k} \in \R^{3 \times 32 \times 32}$ has 3{,}072 dimensions. In such high-dimensional spaces, even non-orthogonal patterns can achieve reliable separation with appropriate non-linear decoders.
\end{enumerate}

The orthogonality regularization loss ($\Lcal_\text{ortho}$) decreasing from 8{,}728 to 11.4 shows that patterns become \emph{approximately} orthogonal, but the Extractor does not require strict orthogonality for reliable decoding.


% ===========================================================================
%  V. DISCUSSION
% ===========================================================================
\section{Discussion}
\label{sec:discussion}

\subsection{Sampling-Path Encoding as a Steganographic Channel}
Our results demonstrate that the output of a frozen diffusion model's sampling chain contains sufficient representational capacity for reliable steganography.
The learned spatial perturbation patterns effectively create a 256-bit channel within the sampling ``path'' without modifying any component of the generative model.
The rapid convergence (BER $\to$ 0 within 440 steps) suggests that the perturbation--extraction mapping is inherently learnable given the architectural alignment between Controller and Extractor.

\subsection{Quality--Robustness Tradeoff}
The empirical results clearly illustrate this tradeoff.
PSNR reaches a peak of 37.6\,dB around step 4{,}000, when only 10\% of batches are subjected to attacks, but declines to 35.0\,dB by step 7{,}000 as the attack probability increases to 40\% and the Controller is compelled to strengthen its perturbation patterns.
During Phase~3, the cosine learning-rate schedule allows the model to recover much of this degradation, stabilizing at 36.7\,dB---approximately 1\,dB below the clean-training peak.

\subsection{STE Proxy Limitations}
Evaluation with real JPEG compression reveals a noticeable gap between the STE-based training proxy and actual codec behavior.
While the proxy transfers effectively at high quality factors ($Q \geq 90$), it consistently underestimates degradation when compression becomes more aggressive ($Q \leq 70$).
Bridging this gap will likely require a more accurate differentiable approximation, such as a compact neural model trained to replicate real codec input--output mappings.

\subsection{Detectability and the Capacity--Undetectability Tradeoff}
The steganalysis classifier achieving AUC = 1.0 is a sobering outcome, though not entirely unexpected for a method operating at this payload level.
Although the tiled perturbations are visually imperceptible, the $8 \times 8$ grid structure leaves a consistent statistical signature that a trained ResNet-18 can readily exploit.
This observation aligns with prior findings that increasing payload capacity often comes at the cost of statistical undetectability~\cite{boroumand2019deep}.
Potential mitigation strategies include: (1)~incorporating an adversarial discriminator during training to encourage more natural perturbation distributions; (2)~adding a loss term that penalizes measurable divergence from clean image statistics; or (3)~reducing per-cell perturbation strength, trading some robustness for improved stealth.

\subsection{FID as a Distributional Metric}
An FID score of 13.7 indicates that, at the distributional level, clean and stego image populations remain statistically separable, even if individual image pairs appear nearly indistinguishable according to PSNR or SSIM.
This finding is consistent with the steganalysis results: the structured perturbations introduced by the Controller subtly shift the overall data distribution in ways detectable by Inception-based features, despite negligible visual differences on a per-image basis.

\subsection{Orthogonality Regularization}
Throughout training, $\Lcal_\text{ortho}$ decreases from 8{,}728 to 11.4, demonstrating that the four learned patterns within each cell progressively approach orthogonal configurations.
This regularization is critical; without it, pattern overlap would reduce per-bit separability, weakening the effective signal and increasing the likelihood of decoding errors.

\subsection{Crop Vulnerability}
The near-random BER observed under crop-and-resize attacks (BER $\sim 0.45$) reflects an architectural limitation rather than a training deficiency.
Cropping alters or removes portions of the spatial grid, breaking the alignment between the Controller’s tiled encoding and the Extractor’s feature representation.
This mirrors well-known weaknesses of spatial-template watermarking under geometric distortions.
We discuss potential remedies in Section~\ref{sec:limitations}.


% ===========================================================================
%  VI. LIMITATIONS AND FUTURE WORK
% ===========================================================================
\section{Limitations and Future Work}
\label{sec:limitations}

\subsection{Crop Robustness}
The spatial tiling design is fundamentally susceptible to crop-and-resize attacks due to spatial misalignment.
Future work may mitigate this limitation through several strategies: (1)~redundant encoding, in which each bit is embedded across multiple non-adjacent cells to provide spatial redundancy; (2)~a global encoding scheme combined with a learned spatial transformer to re-align inputs prior to extraction; or (3)~frequency-domain encoding that distributes information across spatial frequencies rather than anchoring it to fixed grid locations.

\subsection{Real JPEG Robustness}
Although the STE-based JPEG proxy is sufficient for enabling gradient flow during training, empirical evaluation reveals its limitations at quality factors below 70.
A promising direction would be to develop a more faithful neural JPEG surrogate trained on real codec input--output pairs, which could better approximate compression artifacts and improve robustness under aggressive quality settings.

\subsection{Statistical Undetectability}
The fact that a simple classifier can separate clean from stego images with perfect accuracy (AUC = 1.0) is the most pressing practical concern.
Training against a discriminator, or borrowing cover-source mismatch ideas from classical steganography, could bring detectability down, though we expect some cost in either capacity or visual quality.

\subsection{Error-Correcting Codes}
Wrapping the payload in an error-correcting code (Reed-Solomon, LDPC, or similar) would let the system tolerate some bit errors at the cost of reduced effective capacity.
If the ECC layer is made differentiable, it could be folded directly into the end-to-end training loop.

\subsection{Scaling}
All experiments in this paper use $256 \times 256$ images with payloads of 128--512 bits.
Moving to higher resolutions like $512 \times 512$---with a correspondingly larger grid---or porting the approach to latent diffusion models such as Stable Diffusion would open the door to bigger payloads and the image sizes that practical applications demand.


% ===========================================================================
%  VII. ETHICAL CONSIDERATIONS
% ===========================================================================
\section{Ethical Considerations}
\label{sec:ethics}

Like most security tools, steganography sits squarely in dual-use territory.
While it has legitimate applications in digital watermarking, copyright protection, covert communication for journalism and activism, and secure data transmission, it can also be misused for covert exfiltration of sensitive data or evasion of content moderation.

We acknowledge this dual-use nature and advocate for responsible deployment:
\begin{itemize}[itemsep=1pt, topsep=2pt]
    \item We commit to releasing code in a controlled manner, potentially with capacity-limited demonstrations.
    \item We demonstrate that differentiable attack training inherently provides a blueprint for building steganalysis detectors, as the same attack layer can be repurposed for detection training.
    \item We encourage the development of corresponding steganalysis methods and support their publication alongside steganographic advances.
    \item We note that forensic analysis of diffusion-generated images is an active area of research~\cite{corvi2023detection}, and DSCSP's perturbation patterns may be detectable by sufficiently powerful classifiers.
\end{itemize}


% ===========================================================================
%  VIII. CONCLUSION
% ===========================================================================
\section{Conclusion}
\label{sec:conclusion}

We presented DSCSP, a novel steganographic framework that encodes secret messages into the sampling output of frozen pre-trained diffusion models via learned spatial perturbation patterns.
Our multi-bit spatial tiling Controller achieves a 256-bit payload capacity with only $\sim$20K trainable parameters, while a spatially-aligned ResNet-18 Extractor recovers bits with perfect accuracy (BER = 0.0000) at 36.7\,dB PSNR, 0.902 SSIM, 0.032 LPIPS, and FID 13.7.

Through a three-phase robustness curriculum using five differentiable attack proxies, DSCSP learns to resist Gaussian noise (BER = 0.0000 across all tested $\sigma$), blur (BER $< 0.01$ up to radius 1.5), brightness/contrast adjustments (BER $< 0.003$), and real JPEG compression at $Q \geq 90$ (BER $< 0.01$), while identifying crop-and-resize and aggressive JPEG ($Q < 70$) as remaining challenges.
Capacity scaling experiments demonstrate the framework's flexibility from 128 to 512 bits, and an information-theoretic analysis provides insight into how the system exceeds classical channel capacity bounds through learned non-linear decoding.

Steganalysis evaluation reveals a key limitation: stego images are trivially detectable by a trained classifier (AUC = 1.0), highlighting the fundamental tension between payload capacity and statistical undetectability that future work must address.

End-to-end training completes in approximately 100 minutes on a laptop GPU with 6\,GB of memory, making the approach accessible to researchers without extensive computational resources.
We view sampling-path encoding as a practical and extensible framework for diffusion-based steganography, with clear avenues for future development, including adversarial training for improved undetectability, more faithful JPEG proxies, crop-resilient encoding strategies, and extension to higher-resolution image synthesis.


% ===========================================================================
%  REFERENCES
% ===========================================================================
\bibliographystyle{IEEEtran}
\begin{thebibliography}{99}

\bibitem{bengio2013estimating}
Y.~Bengio, N.~L{\'e}onard, and A.~Courville, ``Estimating or propagating gradients through stochastic neurons for conditional computation,'' \emph{arXiv preprint arXiv:1308.3432}, 2013.

\bibitem{boroumand2019deep}
M.~Boroumand, M.~Chen, and J.~Fridrich, ``Deep residual network for steganalysis of digital images,'' \emph{IEEE Trans. Inf. Forensics Security}, vol.~14, no.~5, pp.~1181--1193, 2019.

\bibitem{chan2004hiding}
C.-K.~Chan and L.-M.~Cheng, ``Hiding data in images by simple {LSB} substitution,'' \emph{Pattern Recognition}, vol.~37, no.~3, pp.~469--474, 2004.

\bibitem{corvi2023detection}
R.~Corvi, D.~Cozzolino, G.~Zingarini, G.~Poggi, K.~Nagano, and L.~Verdoliva, ``On the detection of synthetic images generated by diffusion models,'' in \emph{Proc. IEEE Int. Conf. Acoust., Speech Signal Process. (ICASSP)}, 2023, pp.~1--5.

\bibitem{cox2007digital}
I.~Cox, M.~Miller, J.~Bloom, J.~Fridrich, and T.~Kalker, \emph{Digital Watermarking and Steganography}, 2nd~ed.\hskip 1em plus 0.5em minus 0.4em Morgan Kaufmann, 2007.

\bibitem{he2016deep}
K.~He, X.~Zhang, S.~Ren, and J.~Sun, ``Deep residual learning for image recognition,'' in \emph{Proc. IEEE Conf. Comput. Vis. Pattern Recognit. (CVPR)}, 2016, pp.~770--778.

\bibitem{heusel2017gans}
M.~Heusel, H.~Ramsauer, T.~Unterthiner, B.~Nessler, and S.~Hochreiter, ``{GANs} trained by a two time-scale update rule converge to a local {Nash} equilibrium,'' in \emph{Proc. Adv. Neural Inf. Process. Syst. (NeurIPS)}, 2017, pp.~6626--6637.

\bibitem{ho2020denoising}
J.~Ho, A.~Jain, and P.~Abbeel, ``Denoising diffusion probabilistic models,'' in \emph{Proc. Adv. Neural Inf. Process. Syst. (NeurIPS)}, 2020, pp.~6840--6851.

\bibitem{holub2012designing}
V.~Holub and J.~Fridrich, ``Designing steganographic distortion using directional filters,'' in \emph{Proc. IEEE Int. Workshop Inf. Forensics Security (WIFS)}, 2012, pp.~234--239.

\bibitem{holub2014universal}
V.~Holub, J.~Fridrich, and T.~Denemark, ``Universal distortion function for steganography in an arbitrary domain,'' \emph{EURASIP J. Inf. Security}, vol.~2014, no.~1, pp.~1--13, 2014.

\bibitem{kingma2014adam}
D.~P.~Kingma and J.~Ba, ``Adam: A method for stochastic optimization,'' \emph{arXiv preprint arXiv:1412.6980}, 2014.

\bibitem{loshchilov2017sgdr}
I.~Loshchilov and F.~Hutter, ``{SGDR}: Stochastic gradient descent with warm restarts,'' in \emph{Proc. Int. Conf. Learn. Represent. (ICLR)}, 2017.

\bibitem{peng2023stegodiffusion}
W.~Peng, Z.~Li, Y.~Chen, and Z.~Wang, ``Generative steganography diffusion,'' \emph{arXiv preprint arXiv:2305.03472}, 2023.

\bibitem{pevny2010using}
T.~Pevn{\'y}, T.~Filler, and P.~Bas, ``Using high-dimensional image models to perform highly undetectable steganography,'' in \emph{Information Hiding}, Springer, 2010, pp.~161--177.

\bibitem{shin2017jpeg}
R.~Shin and D.~Song, ``{JPEG}-resistant adversarial images,'' in \emph{NeurIPS Workshop Mach. Learn. Comput. Security}, 2017.

\bibitem{song2020denoising}
J.~Song, C.~Meng, and S.~Ermon, ``Denoising diffusion implicit models,'' in \emph{Proc. Int. Conf. Learn. Represent. (ICLR)}, 2021.

\bibitem{wang2004image}
Z.~Wang, A.~C.~Bovik, H.~R.~Sheikh, and E.~P.~Simoncelli, ``Image quality assessment: From error visibility to structural similarity,'' \emph{IEEE Trans. Image Process.}, vol.~13, no.~4, pp.~600--612, 2004.

\bibitem{wei2023diffusionstego}
R.~Wei, C.~Wen, Y.~Xue, and T.~Wang, ``Diffusion-Stego: Training-free diffusion generative steganography via message projection,'' \emph{arXiv preprint arXiv:2305.18726}, 2023.

\bibitem{zhang2018unreasonable}
R.~Zhang, P.~Isola, A.~A.~Efros, E.~Shechtman, and O.~Wang, ``The unreasonable effectiveness of deep features as a perceptual metric,'' in \emph{Proc. IEEE Conf. Comput. Vis. Pattern Recognit. (CVPR)}, 2018, pp.~586--595.

\bibitem{zhang2019steganogan}
K.~A.~Zhang, A.~Cuesta-Infante, L.~Xu, and K.~Veeramachaneni, ``{SteganoGAN}: High capacity image steganography with {GANs},'' \emph{arXiv preprint arXiv:1901.03892}, 2019.

\bibitem{zhang2024ldstega}
J.~Zhang, Y.~Qin, Z.~Liu, and T.~Wang, ``{LDStega}: Practical and robust generative image steganography based on latent diffusion models,'' in \emph{Proc. ACM Multimedia}, 2024.

\bibitem{zhao2024dero}
J.~Zhao, Y.~Liu, and X.~Zhang, ``{DERO}: Diffusion-model-erasure robust watermarking,'' in \emph{Proc. ACM Multimedia}, 2024.

\bibitem{zhu2018hidden}
J.~Zhu, R.~Kaplan, J.~Johnson, and L.~Fei-Fei, ``{HiDDeN}: Hiding data with deep networks,'' in \emph{Proc. Eur. Conf. Comput. Vis. (ECCV)}, 2018, pp.~657--672.

\end{thebibliography}

\end{document}
